\documentclass[12pt,a4paper]{article}
\usepackage[utf8]{inputenc}
\usepackage[T1]{fontenc}
\usepackage[turkish]{babel}
\usepackage{geometry}
\usepackage{enumitem}
\usepackage{hyperref}
\usepackage{titlesec}
\usepackage{fancyhdr}
\usepackage{graphicx}
\usepackage{tabularx}
\usepackage{booktabs}

\geometry{margin=2.5cm}
\pagestyle{fancy}
\fancyhf{}
\rhead{WIND Sentinel}
\lhead{Teknoloji Se\c{c}imi ve Gerek\c{c}elendirme}
\rfoot{\thepage}

\titleformat{\section}{\large\bfseries}{\thesection.}{0.5em}{}
\titleformat{\subsection}{\normalsize\bfseries}{\thesubsection}{0.5em}{}

\begin{document}

\begin{center}
    {\LARGE\bfseries Teknoloji Se\c{c}imi ve Gerek\c{c}elendirme Dok\"uman{\i}}\\[0.5em]
    {\large WIND Sentinel -- R\"uzgar T\"urbini Erken Ar{\i}za Tespit Sistemi}\\[1em]
    \begin{tabular}{ll}
        \textbf{\"O\u{g}renci:} & {\.I}layda Nur U\c{c}ar \\
        \textbf{Tarih:} & 27 Mart 2026 \\
    \end{tabular}
\end{center}

\vspace{1em}
\hrule
\vspace{1em}

\section{Mobil Geli\c{s}tirme Teknolojisi}

\textbf{Se\c{c}im:} Cross-platform -- \textbf{React Native (Expo)}

React Native, tek bir JavaScript/TypeScript kod taban{\i}yla hem Android hem iOS platformlar{\i}na uygulama \"uretmeye olanak tan{\i}r. Expo framework'\"u ise build, test ve da\u{g}{\i}t{\i}m s\"ure\c{c}lerini basitle\c{s}tirir.

\section{Veri Kayna\u{g}{\i}}

\textbf{Se\c{c}im:} Bulut tabanl{\i} veritaban{\i} -- \textbf{PostgreSQL (Docker container)}

Uygulama verilerini merkezi bir PostgreSQL veritaban{\i}nda saklar. Mobil uygulama veritaban{\i}na do\u{g}rudan eri\c{s}mez; t\"um veri al{\i}\c{s}veri\c{s}i \textbf{API Gateway} (Node.js/Express) \"uzerinden RESTful API ve WebSocket ba\u{g}lant{\i}lar{\i}yla ger\c{c}ekle\c{s}ir. Bu yap{\i}, veri g\"uvenli\u{g}ini ve tutarl{\i}l{\i}\u{g}{\i}n{\i} sa\u{g}lar.

\section{Ek Ara\c{c}lar ve K\"ut\"uphaneler}

\begin{table}[h]
\centering
\renewcommand{\arraystretch}{1.3}
\begin{tabularx}{\textwidth}{lX}
\toprule
\textbf{Ara\c{c} / K\"ut\"uphane} & \textbf{Kullan{\i}m Amac{\i}} \\
\midrule
React Navigation & Ekranlar aras{\i} ge\c{c}i\c{s} ve sekme y\"onetimi \\
Socket.IO Client & WebSocket ile ger\c{c}ek zamanl{\i} alarm bildirimleri \\
Expo Notifications & Push bildirim altyap{\i}s{\i} \\
AsyncStorage & JWT token ve kullan{\i}c{\i} tercihlerinin yerel saklanmas{\i} \\
Axios & API Gateway ile HTTP ileti\c{s}imi \\
Lucide Icons & Tutarl{\i} ve modern ikon seti \\
Docker \& Docker Compose & T\"um mikroservislerin konteynerle\c{s}tirilmesi \\
RabbitMQ & Servisler aras{\i} asenkron mesajla\c{s}ma \\
Isolation Forest + XGBoost & Hibrit anomali tespit modeli \\
\bottomrule
\end{tabularx}
\end{table}

\section{Gerek\c{c}elendirme}

\subsection{Neden React Native (Expo)?}

\begin{itemize}[leftmargin=1.5em, itemsep=0.3em]
    \item \textbf{Problemin yap{\i}s{\i}:} Uygulama; dashboard g\"or\"unt\"uleme, alarm listeleme, t\"urbin durumu izleme ve manuel \"ol\c{c}\"um giri\c{s}i gibi aray\"uz a\u{g}{\i}rl{\i}kl{\i} i\c{s}levler i\c{c}erir. Kamera, Bluetooth veya cihaz d\"uzeyinde d\"u\c{s}\"uk seviyeli eri\c{s}im gerektirmez. Bu nedenle native geli\c{s}tirme gereksiz karma\c{s}{\i}kl{\i}k ekler.
    \item \textbf{Tek kod taban{\i}:} Web dashboard zaten React ile geli\c{s}tirildi\u{g}inden, React Native'e ge\c{c}i\c{s} s\"ures\"uz yak{\i}nd{\i}r. Bile\c{s}en yap{\i}s{\i}, state y\"onetimi ve API servis katman{\i} do\u{g}rudan yeniden kullan{\i}labilir.
    \item \textbf{Geli\c{s}tirme s\"uresi:} Expo'nun sa\u{g}lad{\i}\u{g}{\i} haz{\i}r build pipeline, hot reload ve OTA (Over-The-Air) g\"uncelleme deste\u{g}i, tek ki\c{s}ilik geli\c{s}tirme ekibinde s\"ureyi ciddi \c{s}ekilde k{\i}salt{\i}r.
    \item \textbf{Kullan{\i}c{\i} say{\i}s{\i} beklentisi:} Sistem bir r\"uzgar \c{c}iftli\u{g}inin bak{\i}m ekibine (5--20 teknisyen) y\"oneliktir. Y\"uksek performans optimizasyonu gerektiren bir kullan{\i}c{\i} taban{\i} s\"oz konusu de\u{g}ildir.
\end{itemize}

\subsection{Neden Bulut Tabanl{\i} Veritaban{\i}?}

\begin{itemize}[leftmargin=1.5em, itemsep=0.3em]
    \item \textbf{Ger\c{c}ek zamanl{\i} veri:} T\"urbin sens\"orleri her 10 dakikada veri \"uretir. Bu verinin i\c{s}lenmesi, analiz edilmesi ve sonu\c{c}lar{\i}n hem web hem mobil istemcilere anl{\i}k iletilmesi gerekir. Yerel veritaban{\i} bu senkronizasyonu sa\u{g}layamaz.
    \item \textbf{Veri b\"ut\"unl\"u\u{g}\"u:} Birden fazla kullan{\i}c{\i} ayn{\i} alarm{\i} g\"or\"up i\c{s}lem yapabilir. Merkezi veritaban{\i}, foreign key, unique constraint ve check constraint ile veri tutarl{\i}l{\i}\u{g}{\i}n{\i} garanti eder.
    \item \textbf{Mikroservis mimarisi:} Sistemin 5 ba\u{g}{\i}ms{\i}z servisi (Ingestion, Feature, Prediction, Notification, API Gateway) ayn{\i} veritaban{\i}na eri\c{s}ir. Bu yap{\i} ancak merkezi bir veritaban{\i}yla m\"umk\"und\"ur.
\end{itemize}

\subsection{Neden RabbitMQ?}

\begin{itemize}[leftmargin=1.5em, itemsep=0.3em]
    \item \textbf{Asenkron i\c{s}leme:} Sens\"or verisi i\c{s}leme pipeline'{\i} (ham veri $\rightarrow$ \"ozellik \c{c}{\i}kar{\i}m{\i} $\rightarrow$ tahmin $\rightarrow$ bildirim) s{\i}ral{\i} ama ba\u{g}{\i}ms{\i}z ad{\i}mlardan olu\c{s}ur. RabbitMQ, bir servis yava\c{s}lasa bile mesajlar{\i}n kuyrukta bekledi\u{g}i, veri kayb{\i} ya\c{s}anmad{\i}\u{g}{\i} bir yap{\i} sunar.
    \item \textbf{Dayan{\i}kl{\i}l{\i}k:} Durable queue ve dead letter exchange mekanizmalar{\i} sayesinde hatal{\i} mesajlar kaybolmaz, incelenmek \"uzere ayr{\i} kuyru\u{g}a y\"onlendirilir.
\end{itemize}

\end{document}
